\nonstopmode{}
\documentclass[a4paper]{book}
\usepackage[times,inconsolata,hyper]{Rd}
\usepackage{makeidx}
\makeatletter\@ifl@t@r\fmtversion{2018/04/01}{}{\usepackage[utf8]{inputenc}}\makeatother
% \usepackage{graphicx} % @USE GRAPHICX@
\makeindex{}
\begin{document}
\chapter*{}
\begin{center}
{\textbf{\huge Package `qtbi'}}
\par\bigskip{\large \today}
\end{center}
\ifthenelse{\boolean{Rd@use@hyper}}{\hypersetup{pdftitle = {qtbi: Quantum Toxic Burden Index}}}{}
\ifthenelse{\boolean{Rd@use@hyper}}{\hypersetup{pdfauthor = {January G. Msemakweli}}}{}
\begin{description}
\raggedright{}
\item[Title]\AsIs{Quantum Toxic Burden Index}
\item[Version]\AsIs{0.1.2}
\item[Description]\AsIs{Compute the Quantum Toxic Burden Index from multi-exposure
panels using a fixed quantum-inspired entanglement encoder; method
reference (2026) <}\Rhref{https://doi.org/10.5281/zenodo.20476574}{doi:10.5281/zenodo.20476574}\AsIs{>. Provides percentile
encoding, optional potency-weighted readout, and synergy diagnostics
for environmental mixture burden scores.}
\item[License]\AsIs{MIT + file LICENSE}
\item[URL]\AsIs{}\url{https://github.com/january-msemakweli/qtbi}\AsIs{,
}\url{https://doi.org/10.5281/zenodo.20476574}\AsIs{}
\item[BugReports]\AsIs{}\url{https://github.com/january-msemakweli/qtbi/issues}\AsIs{}
\item[Encoding]\AsIs{UTF-8}
\item[Language]\AsIs{en-US}
\item[Roxygen]\AsIs{list(markdown = TRUE)}
\item[RoxygenNote]\AsIs{7.3.3}
\item[Depends]\AsIs{R (>= 4.1.0)}
\item[Suggests]\AsIs{testthat (>= 3.0.0)}
\item[Config/testthat/edition]\AsIs{3}
\item[NeedsCompilation]\AsIs{no}
\item[Author]\AsIs{January G. Msemakweli [aut, cre]}
\item[Maintainer]\AsIs{January G. Msemakweli }\email{jmsemak1@jh.edu}\AsIs{}
\end{description}
\Rdcontents{Contents}
\HeaderA{qtbi-package}{Quantum Toxic Burden Index}{qtbi.Rdash.package}
\aliasA{qtbi}{qtbi-package}{qtbi}
\keyword{internal}{qtbi-package}
%
\begin{Description}
Compute the Quantum Toxic Burden Index (QTBI) from multi-exposure panels
using a fixed quantum-inspired entanglement encoder.
\end{Description}
%
\begin{Section}{Typical workflow}

\begin{itemize}

\item{} \code{\LinkA{estimate\_qtbi()}{estimate.Rul.qtbi}} adds percentile and QTBI columns to your data frame.
\item{} Optional \code{reference\_doses} enable potency-weighted readout at the index step.
Weights are rescaled automatically so weighted QTBI stays on the same
\AsIs{\texttt{[0, n]}} scale as the unweighted index.
\item{} \code{\LinkA{diagnose\_qtbi()}{diagnose.Rul.qtbi}} summarizes synergy sensitivity and monotonicity.

\end{itemize}

\end{Section}
%
\begin{Author}
\strong{Maintainer}: January G. Msemakweli \email{jmsemak1@jh.edu}

\end{Author}
%
\begin{SeeAlso}
Useful links:
\begin{itemize}

\item{} \url{https://github.com/january-msemakweli/qtbi}
\item{} \Rhref{https://doi.org/10.5281/zenodo.20476574}{doi:10.5281\slash{}zenodo.20476574}
\item{} Report bugs at \url{https://github.com/january-msemakweli/qtbi/issues}

\end{itemize}


\end{SeeAlso}
\HeaderA{build\_statevector}{Build the n-qubit entanglement statevector for one subject}{build.Rul.statevector}
%
\begin{Description}
Build the n-qubit entanglement statevector for one subject
\end{Description}
%
\begin{Usage}
\begin{verbatim}
build_statevector(percentiles, synergy = 0.6)
\end{verbatim}
\end{Usage}
%
\begin{Arguments}
\begin{ldescription}
\item[\code{percentiles}] Numeric vector of exposure percentiles in \AsIs{\texttt{[0, 1]}}.

\item[\code{synergy}] Synergy strength in \AsIs{\texttt{[0, 1]}}.
\end{ldescription}
\end{Arguments}
%
\begin{Value}
Complex statevector of length \code{2\textasciicircum{}n}.
\end{Value}
\HeaderA{circuit\_gate\_schedule}{Gate schedule for circuit diagrams}{circuit.Rul.gate.Rul.schedule}
%
\begin{Description}
Gate schedule for circuit diagrams
\end{Description}
%
\begin{Usage}
\begin{verbatim}
circuit_gate_schedule(n_exposures = 4L, n_metals = NULL)
\end{verbatim}
\end{Usage}
%
\begin{Arguments}
\begin{ldescription}
\item[\code{n\_exposures}] Number of exposures (qubits).

\item[\code{n\_metals}] Deprecated alias for \code{n\_exposures}.
\end{ldescription}
\end{Arguments}
%
\begin{Value}
List with gate schedule metadata.
\end{Value}
\HeaderA{diagnose\_qtbi}{Diagnose QTBI encoder behavior on processed data}{diagnose.Rul.qtbi}
%
\begin{Description}
Diagnose QTBI encoder behavior on processed data
\end{Description}
%
\begin{Usage}
\begin{verbatim}
diagnose_qtbi(
  data,
  synergy_grid = seq(0, 1, by = 0.05),
  synergy_ref = NULL,
  verbose = TRUE
)
\end{verbatim}
\end{Usage}
%
\begin{Arguments}
\begin{ldescription}
\item[\code{data}] A \code{qtbi\_data} object from \code{\LinkA{estimate\_qtbi()}{estimate.Rul.qtbi}}.

\item[\code{synergy\_grid}] Grid of synergy values for sensitivity curves.

\item[\code{synergy\_ref}] Reference synergy for summary statistics (defaults to
the value used in \code{\LinkA{estimate\_qtbi()}{estimate.Rul.qtbi}}).

\item[\code{verbose}] If \code{TRUE}, print a concise summary.
\end{ldescription}
\end{Arguments}
%
\begin{Value}
A \code{qtbi\_diagnosis} object (also printed when \code{verbose = TRUE}).
\end{Value}
\HeaderA{estimate\_qtbi}{Estimate QTBI and append scores to a data frame}{estimate.Rul.qtbi}
%
\begin{Description}
Computes within-cohort exposure percentiles, runs the fixed entanglement
encoder, and returns the input data with percentile and QTBI columns added.
\end{Description}
%
\begin{Usage}
\begin{verbatim}
estimate_qtbi(
  data,
  chemicals,
  synergy_strength = 0.6,
  qtbi_col = "qtbi",
  percentile_prefix = "pct_",
  exposure_names = NULL,
  pct_cols = NULL,
  reference_doses = NULL,
  reference_index = NULL
)
\end{verbatim}
\end{Usage}
%
\begin{Arguments}
\begin{ldescription}
\item[\code{data}] A data frame containing exposure columns.

\item[\code{chemicals}] Character vector of exposure column names.

\item[\code{synergy\_strength}] Synergy strength in \AsIs{\texttt{[0, 1]}}.

\item[\code{qtbi\_col}] Name for the QTBI score column (default \code{"qtbi"}).

\item[\code{percentile\_prefix}] Prefix for percentile columns (default \code{"pct\_"}).

\item[\code{exposure\_names}] Optional display names for plots (defaults to \code{chemicals}).

\item[\code{pct\_cols}] Optional percentile column names (same length as \code{chemicals}).

\item[\code{reference\_doses}] Optional named numeric vector of oral reference doses
in mg/kg/day, with names matching \code{exposure\_names} (or \code{chemicals} if
\code{exposure\_names} is omitted). The package derives potency weights as
\code{reference\_index} dose divided by each component dose, rescales them so
their sum equals the number of exposures (keeping QTBI on the same
\eqn{[0, n]}{} scale as the unweighted index), and applies them at readout only.

\item[\code{reference\_index}] Name of the index chemical for potency ratios. Defaults
to the first \code{exposure\_names} entry (or first \code{chemicals} entry).
\end{ldescription}
\end{Arguments}
%
\begin{Value}
A \code{qtbi\_data} object (data frame) with QTBI and percentile columns.
\end{Value}
%
\begin{Examples}
\begin{ExampleCode}
df <- data.frame(
  Pb = c(1, 2, 3, 4),
  As = c(4, 3, 2, 1),
  Cd = c(2, 2, 3, 3),
  Hg = c(1, 3, 2, 4)
)
out <- estimate_qtbi(
  df,
  chemicals = c("Pb", "As", "Cd", "Hg"),
  synergy_strength = 0.6,
  reference_doses = c(Pb = 6.3e-4, As = 6.0e-5, Cd = 5.0e-4, Hg = 1.0e-4),
  reference_index = "Pb"
)
out$qtbi
\end{ExampleCode}
\end{Examples}
\HeaderA{exposure\_percentile}{Convert exposures to within-cohort percentiles}{exposure.Rul.percentile}
%
\begin{Description}
Convert exposures to within-cohort percentiles
\end{Description}
%
\begin{Usage}
\begin{verbatim}
exposure_percentile(x)
\end{verbatim}
\end{Usage}
%
\begin{Arguments}
\begin{ldescription}
\item[\code{x}] Numeric exposure vector.
\end{ldescription}
\end{Arguments}
%
\begin{Value}
Numeric percentiles in \AsIs{\texttt{(0, 1]}}.
\end{Value}
\HeaderA{is\_qtbi\_data}{Test for a qtbi\_data object}{is.Rul.qtbi.Rul.data}
%
\begin{Description}
Test for a qtbi\_data object
\end{Description}
%
\begin{Usage}
\begin{verbatim}
is_qtbi_data(x)
\end{verbatim}
\end{Usage}
%
\begin{Arguments}
\begin{ldescription}
\item[\code{x}] Object to test.
\end{ldescription}
\end{Arguments}
%
\begin{Value}
\code{TRUE} if \code{x} inherits from \code{"qtbi\_data"}.
\end{Value}
\HeaderA{marginal\_toxic\_probs}{Marginal toxic probabilities from a QTBI statevector}{marginal.Rul.toxic.Rul.probs}
%
\begin{Description}
Marginal toxic probabilities from a QTBI statevector
\end{Description}
%
\begin{Usage}
\begin{verbatim}
marginal_toxic_probs(state)
\end{verbatim}
\end{Usage}
%
\begin{Arguments}
\begin{ldescription}
\item[\code{state}] Complex statevector from \code{\LinkA{build\_statevector()}{build.Rul.statevector}}.
\end{ldescription}
\end{Arguments}
%
\begin{Value}
Numeric vector of marginal toxic probabilities.
\end{Value}
\HeaderA{normalize\_potency\_weights}{Rescale potency weights to match the unweighted QTBI scale}{normalize.Rul.potency.Rul.weights}
%
\begin{Description}
When potency weights are applied at readout, rescale them so their sum equals
the number of exposures in the panel. This keeps weighted QTBI on the same
\eqn{[0, n]}{} scale as the unweighted sum of marginals, which aids comparison
between weighted and unweighted analyses.
\end{Description}
%
\begin{Usage}
\begin{verbatim}
normalize_potency_weights(weights, target_sum = NULL)
\end{verbatim}
\end{Usage}
%
\begin{Arguments}
\begin{ldescription}
\item[\code{weights}] Named numeric vector of raw potency ratios from
\code{\LinkA{potency\_weights\_from\_reference\_doses()}{potency.Rul.weights.Rul.from.Rul.reference.Rul.doses}}.

\item[\code{target\_sum}] Target sum for the rescaled weights. Defaults to the
number of exposures (\code{length(weights)}).
\end{ldescription}
\end{Arguments}
%
\begin{Value}
Named numeric vector of rescaled weights aligned with \code{weights}.
\end{Value}
\HeaderA{percentile\_matrix}{Build percentile matrix from a data frame}{percentile.Rul.matrix}
%
\begin{Description}
Build percentile matrix from a data frame
\end{Description}
%
\begin{Usage}
\begin{verbatim}
percentile_matrix(data, chemicals)
\end{verbatim}
\end{Usage}
%
\begin{Arguments}
\begin{ldescription}
\item[\code{data}] Data frame containing exposure columns.

\item[\code{chemicals}] Character vector of column names.
\end{ldescription}
\end{Arguments}
%
\begin{Value}
Numeric matrix with one percentile column per chemical.
\end{Value}
\HeaderA{potency\_weights\_from\_reference\_doses}{Derive relative potency weights from reference doses}{potency.Rul.weights.Rul.from.Rul.reference.Rul.doses}
%
\begin{Description}
Computes outcome-independent weights as the ratio of the index chemical
reference dose to each component reference dose:
\deqn{w_i = \mathrm{RfD}_{\mathrm{index}} / \mathrm{RfD}_i}{}
\end{Description}
%
\begin{Usage}
\begin{verbatim}
potency_weights_from_reference_doses(
  reference_doses,
  chemicals,
  reference_index = NULL
)
\end{verbatim}
\end{Usage}
%
\begin{Arguments}
\begin{ldescription}
\item[\code{reference\_doses}] Named numeric vector of oral reference doses in
mg/kg/day, with one value per exposure in \code{chemicals}.

\item[\code{chemicals}] Character vector of exposure identifiers in panel order.
Names must match \code{names(reference\_doses)}.

\item[\code{reference\_index}] Name of the index chemical. Defaults to the first
entry in \code{chemicals}.
\end{ldescription}
\end{Arguments}
%
\begin{Details}
Lower reference doses (more potent toxicants) receive larger weights.
\end{Details}
%
\begin{Value}
Named numeric vector of potency weights aligned with \code{chemicals}.
\end{Value}
\HeaderA{qtbi\_from\_pcts}{Compute QTBI from percentile matrix}{qtbi.Rul.from.Rul.pcts}
%
\begin{Description}
Compute QTBI from percentile matrix
\end{Description}
%
\begin{Usage}
\begin{verbatim}
qtbi_from_pcts(pct_mat, synergy = 0.6, weights = NULL)
\end{verbatim}
\end{Usage}
%
\begin{Arguments}
\begin{ldescription}
\item[\code{pct\_mat}] Numeric matrix of within-cohort percentiles (rows = subjects).

\item[\code{synergy}] Synergy strength in \AsIs{\texttt{[0, 1]}}.

\item[\code{weights}] Optional named numeric potency weights aligned with columns.
\end{ldescription}
\end{Arguments}
%
\begin{Value}
Numeric vector of QTBI scores.
\end{Value}
\HeaderA{qtbi\_from\_state}{Compute QTBI from a statevector}{qtbi.Rul.from.Rul.state}
%
\begin{Description}
Compute QTBI from a statevector
\end{Description}
%
\begin{Usage}
\begin{verbatim}
qtbi_from_state(state, weights = NULL)
\end{verbatim}
\end{Usage}
%
\begin{Arguments}
\begin{ldescription}
\item[\code{state}] Complex statevector from \code{\LinkA{build\_statevector()}{build.Rul.statevector}}.

\item[\code{weights}] Optional named numeric potency weights aligned with qubits.
When \code{NULL}, marginals are summed with equal weight.
\end{ldescription}
\end{Arguments}
%
\begin{Value}
Scalar QTBI score (sum of marginal toxic probabilities).
\end{Value}
\HeaderA{qtbi\_from\_vector}{Compute QTBI from one percentile vector}{qtbi.Rul.from.Rul.vector}
%
\begin{Description}
Compute QTBI from one percentile vector
\end{Description}
%
\begin{Usage}
\begin{verbatim}
qtbi_from_vector(pct_vec, synergy = 0.6, weights = NULL)
\end{verbatim}
\end{Usage}
%
\begin{Arguments}
\begin{ldescription}
\item[\code{pct\_vec}] Numeric vector of within-cohort percentiles.

\item[\code{synergy}] Synergy strength in \AsIs{\texttt{[0, 1]}}.

\item[\code{weights}] Optional named numeric potency weights aligned with \code{pct\_vec}.
\end{ldescription}
\end{Arguments}
%
\begin{Value}
Scalar QTBI score.
\end{Value}
\HeaderA{qtbi\_help}{Package help index for qtbi}{qtbi.Rul.help}
%
\begin{Description}
Prints a topic index for the main \strong{qtbi} functions and optionally opens
documentation for a specific topic. Use this when you want a guided entry
point to the package; standard R help remains available via \code{?function\_name}
and \code{help(package = "qtbi")}.
\end{Description}
%
\begin{Usage}
\begin{verbatim}
qtbi_help(topic = NULL)
\end{verbatim}
\end{Usage}
%
\begin{Arguments}
\begin{ldescription}
\item[\code{topic}] Optional character string naming a function or topic to open.
If omitted, a summary index is printed.
\end{ldescription}
\end{Arguments}
%
\begin{Value}
Invisibly returns \code{NULL}. When \code{topic} is supplied, the corresponding
help page is opened as a side effect.
\end{Value}
%
\begin{Examples}
\begin{ExampleCode}
qtbi_help()
qtbi_help("estimate_qtbi")
\end{ExampleCode}
\end{Examples}
\HeaderA{qtbi\_meta}{Retrieve QTBI metadata stored on a processed data frame}{qtbi.Rul.meta}
%
\begin{Description}
Retrieve QTBI metadata stored on a processed data frame
\end{Description}
%
\begin{Usage}
\begin{verbatim}
qtbi_meta(data)
\end{verbatim}
\end{Usage}
%
\begin{Arguments}
\begin{ldescription}
\item[\code{data}] A \code{qtbi\_data} object from \code{\LinkA{estimate\_qtbi()}{estimate.Rul.qtbi}}.
\end{ldescription}
\end{Arguments}
%
\begin{Value}
List with chemicals, synergy strength, and percentile column names.
\end{Value}
\HeaderA{synergy\_diagnostics}{Full circuit diagnostics for a cohort percentile matrix}{synergy.Rul.diagnostics}
%
\begin{Description}
Full circuit diagnostics for a cohort percentile matrix
\end{Description}
%
\begin{Usage}
\begin{verbatim}
synergy_diagnostics(
  pct_mat,
  synergy_grid = seq(0, 1, by = 0.05),
  synergy_ref = 0.6,
  exposure_names = NULL,
  metal_names = NULL,
  weights = NULL
)
\end{verbatim}
\end{Usage}
%
\begin{Arguments}
\begin{ldescription}
\item[\code{pct\_mat}] Numeric matrix of within-cohort percentiles (rows = subjects).

\item[\code{synergy\_grid}] Grid of synergy values in \AsIs{\texttt{[0, 1]}}.

\item[\code{synergy\_ref}] Reference synergy for summary statistics.

\item[\code{exposure\_names}] Labels for exposures (columns of \code{pct\_mat}).

\item[\code{metal\_names}] Deprecated alias for \code{exposure\_names}.

\item[\code{weights}] Optional potency weights aligned with columns of \code{pct\_mat}.
\end{ldescription}
\end{Arguments}
%
\begin{Value}
List with cohort\_band, marginals, monotonicity, synergy\_ref, sens\_ref.
\end{Value}
\HeaderA{synergy\_sensitivity}{Synergy sensitivity on cohort median percentiles}{synergy.Rul.sensitivity}
%
\begin{Description}
Synergy sensitivity on cohort median percentiles
\end{Description}
%
\begin{Usage}
\begin{verbatim}
synergy_sensitivity(
  pct_mat,
  synergy_grid = seq(0, 1, by = 0.05),
  weights = NULL
)
\end{verbatim}
\end{Usage}
%
\begin{Arguments}
\begin{ldescription}
\item[\code{pct\_mat}] Numeric matrix of within-cohort percentiles.

\item[\code{synergy\_grid}] Grid of synergy values in \AsIs{\texttt{[0, 1]}}.

\item[\code{weights}] Optional potency weights aligned with columns of \code{pct\_mat}.
\end{ldescription}
\end{Arguments}
%
\begin{Value}
Data frame with synergy, qtbi, and additive\_baseline columns.
\end{Value}
\printindex{}
\end{document}
